\section{Study Design}

\subsection{Overall Design}

This is a prospective, multicenter, randomized, parallel-group, controlled,
open-label study with blinded independent central review (BICR), conducted across
eight academic high-volume hepatopancreatobiliary (HPB) centers. A total of
$n=220$ participants are randomized 1:1 (110 per arm). Arm A receives perioperative
daraxonrasib (RMC-6236) at the recommended Phase 2 dose (RP2D) of 300~mg orally
once daily plus the on-premises large language model (LLM) directed eight-arm
robotic pancreaticoduodenectomy (Whipple) system (PancreSpeed II). Arm B receives
modified FOLFIRINOX \cite{Conroy2018} plus institutional-standard high-volume
pancreaticoduodenectomy that is not LLM-directed. Allocation uses central
permuted-block randomization 1:1, stratified by resectability (resectable versus
borderline-resectable), KRAS allele (G12D versus other G12), neoadjuvant therapy
(yes versus no), and site. The study is conducted under a combined IND/IDE
submission with a single integrated protocol, a single informed consent carrying
the Physical AI opt-out, and a single multicenter safety governance structure.

The device operates at Stage 2 supervised semiautonomy as a Class II collaborative
system under continuous human oversight with an immediate manual override; full
autonomy is prohibited under 21 CFR \S312.21(e), and any expanded
supervised-autonomous subtask requires independent safety review and a Unified
Safety Level (USL) of at least 8.0. A mandatory upgraded Phase 0
simulation-validation gate must be passed at each site before any patient-facing
robotic activity (\S4.3). The eight-site, powered, equitable design is enabled by
the Patient Access and Equity Fund (the Patient-Aligned Co-Investment Facility,
PACIF) held behind a capital firewall, so that funders cannot touch randomization,
endpoint definition, outcome adjudication, or analysis \cite{kawchak2026hr9510,ecfr2026part54}.
Figure 8 presents the randomization and multicenter design.

\subsection{Scientific Rationale for Study Design}

A randomized, controlled, parallel-group design is now the scientifically and
ethically appropriate next step, inverting the logic of the single-arm predicate.
The first-in-human study was single-arm precisely because the early safety profile,
the RP2D, and the technical feasibility of the combined drug-device intervention
were unknown, and randomizing a vulnerable population before that evidence existed
could not be justified. The Phase 1 predicate \cite{kawchak2026phase1} has now
established all three, so genuine clinical equipoise exists between the
investigational regimen and institutional-standard care, and a randomized
controlled comparison is the appropriate and ethical way to test comparative
efficacy. The trial follows the CONSORT framework for reporting a parallel-group
randomized design \cite{Schulz2010consort}.

Open-label conduct is unavoidable because a surgical-system intervention cannot be
masked from the operating team or the participant. Bias control is therefore
concentrated downstream of the intervention: the primary PFS endpoint is supported
by BICR of imaging, the R0 and MPR endpoints are read on central masked pathology,
and all confirmatory endpoints pass through blinded endpoint adjudication, so that
the open-label delivery does not propagate into outcome ascertainment. The internal
control, Arm B, replaces the external Dutch benchmark used in the predicate as the
primary comparator \cite{DutchCohort2025}; the Dutch cohort is retained only as an
external descriptive reference. The study population remains vulnerable, with a
total five-year survival below 13\% \cite{Siegel2025}, so the Patient Access and
Equity Fund and per-cohort DSMB oversight are maintained to protect participants
throughout randomized comparison.

\begin{mermaidfig}
\node[mmin]  (elig) at (0,-0.4)        {Eligible\\participant};
\node[mmstep](rand) at (0,-2.8)     {Central stratified\\randomization 1:1\\(resectability, KRAS\\allele, neoadjuvant, site)};
\node[mmgoal](armA) at (-3.9,-5.2)  {Arm A\\daraxonrasib RP2D +\\LLM robotic Whipple};
\node[mmstep](armB) at (3.9,-5.2)   {Arm B\\mFOLFIRINOX +\\standard Whipple};
\node[mmdark](sites)at (0,-7.5)      {8 academic HPB sites\\(harmonized fleet)};
\node[mmstep](bicr) at (0,-9.2)     {BICR + central\\masked pathology};
\node[mmgoal](adj)  at (0,-11.1)     {Endpoint\\adjudication};
\draw[mmedge]  (elig) -- (rand);
\draw[mmedge]  (rand) -- (armA);
\draw[mmedge]  (rand) -- (armB);
\draw[mmedge]  (armA) -- (sites);
\draw[mmedge]  (armB) -- (sites);
\draw[mmedge]  (sites) -- (bicr);
\draw[mmedge]  (bicr) -- (adj);
\end{mermaidfig}
\vspace{-0.7cm}
\figcaption{Figure 8. Randomization and multicenter design. Eligible participants
undergo central permuted-block randomization 1:1, stratified by resectability,
KRAS allele, neoadjuvant therapy, and site, to Arm A (daraxonrasib at the RP2D plus
the LLM-directed robotic Whipple) or Arm B (modified FOLFIRINOX plus standard
high-volume Whipple). Both arms are delivered across the academic HPB sites
operating as a harmonized fleet, and all outcomes flow through blinded independent
central review and central masked pathology into blinded endpoint adjudication.}

\subsection{Justification for Dose and Device Readiness}

\paragraph{Daraxonrasib at the established RP2D.}
Daraxonrasib is administered at the RP2D of 300~mg orally once daily established in
the Phase 1 predicate and anchored to the RASolute 302 pan-KRAS program
\cite{kawchak2026phase1,rasolute302}. Because the RP2D is already established, there
is no 3+3 dose escalation in this Phase 2 study; every participant in Arm A receives
the same fixed dose, and the trial tests comparative efficacy rather than
dose-finding.

\paragraph{Device ``dose'' analogue and the upgraded Phase 0 gate.}
The device has no pharmacologic dose; its dose analogue is the readiness evidence
required before any patient-facing activity, escalated only by passing the mandatory
upgraded Phase 0 simulation-validation gate. Relative to the predicate, the gate is
upgraded to require a USL rating of at least 8.0, at least 5000 simulated procedures
reproduced across at least three independent simulation frameworks, and a tightened
sim-to-real fidelity of a trajectory gap below 1.5~mm and a tip-force gap below
0.4~N, together with per-site qualification and multicenter fleet harmonization so
that every device instance behaves identically across the eight centers. These
preconditions operate with the intra-operative hard caps (per-arm tip-force caps of
$\leq 3$~N, a cumulative cross-arm cap of $\leq 18$~N, a cross-arm emergency-stop
budget of $\leq 3$~ms, and a five-vessel no-fly gate around the major peripancreatic
vessels) \cite{ecfr2026part812,kawchak2026adapt312}. The expanded Phase 0 compute
required to meet these thresholds is funded by the co-investment facility behind the
capital firewall, which purchases operational capacity only and never scientific
outcomes. Figure 9 presents the staged-autonomy model for Phase 2.

\begin{mermaidfig}
\node[mmin]  (st1) at (-0.05,0)  {Stage 1\\Clinician-controlled\\teleoperation /\\supervised;\\full autonomy\\prohibited \S312.21(e)};
\node[mmdec] (r1)  at (5.1,0)    {Independent safety\\review + USL\\$\geq 8.0$};
\node[mmstep](st2) at (10.35,0)  {Stage 2\\Supervised\\semiautonomous;\\Class II, continuous\\oversight, immediate\\manual override};
\node[mmdec] (r2)  at (10.35,-3.4){Expanded subtask\\independent review +\\USL $\geq 8.0$};
\node[mmgoal](st3) at (5.1,-3.4) {Stage 3\\Higher autonomy\\(reserved, prohibited\\this phase)};
\draw[mmedge] (st1) -- (r1);
\draw[mmedge] (r1) -- (st2);
\draw[mmedge] (st2) -- (r2);
\draw[mmedge] (r2) -- (st3);
\end{mermaidfig}
\vspace{-0.7cm}
\figcaption{Figure 9. Phase-graduated staged-autonomy model for Phase 2. The
operative system runs at Stage 2\\supervised semiautonomy (Class II, continuous
oversight with immediate manual override). Any expanded supervised-autonomous
subtask requires an independent safety review and a Unified Safety Level of at least\\8.0; full autonomy (Stage 3) remains reserved and prohibited at this phase under 21
CFR \S312.21(e).}

\subsection{End of Study Definition}

An individual participant is considered to have completed the study at the last
visit specified in the Schedule of Activities (\S1.3), which is the 24-month
overall-survival assessment, or at the participant's death, documented disease
progression precluding further follow-up, or withdrawal of consent, whichever
occurs first. The end of the study as a whole is defined as the date of the last
scheduled PFS or OS assessment for the last enrolled participant together with the
primary-analysis database lock, which occurs after approximately 140 PFS events
have accrued. After that date, no further study-specific procedures are performed.
Participants who discontinue the investigational intervention for a protocol-defined
reason (\S7) remain in long-term follow-up for survival to the 24-month endpoint
unless they withdraw consent for follow-up.
